\documentclass[11pt]{article}
\usepackage[utf8]{inputenc}
\usepackage[T1]{fontenc}
\usepackage{lmodern}
\usepackage{textcomp}
\usepackage{amsmath, amssymb, amsthm}
\usepackage{graphicx}
\usepackage{listings}
\usepackage[hidelinks]{hyperref}

\theoremstyle{plain}
\newtheorem{theorem}{Teorema}[section]
\newtheorem{lemma}[theorem]{Lema}
\newtheorem{proposition}[theorem]{Proposición}
\newtheorem{corollary}[theorem]{Corolario}

\theoremstyle{definition}
\newtheorem{definition}[theorem]{Definición}


\title{c5}
\date{}
\begin{document}
\maketitle

\section{Chapter 5}

\section{Describing Functions}

We know only a handful of words to describe properties of functions:

injective, surjective, invertible (bijective), constant.

In this chapter we consider attributes of real functions f : R $\to$R. We expand our vocabulary substantially, and introduce important terms that are found in more general settings (periodic, bounded, continuous). Then we export this terminology to the real sequences, which are functions defined over N.

\section{5.1 Ordering Properties}

The real line R is ordered, meaning that for any pair (x, y) of real numbers, precisely one of the three relational expressions

x < y x = y x > y

is true, and the other two are false. We begin with properties that are formulated in terms of ordering.

A simple attribute of a real function is the sign of the values it assumes.

$\forall$x $\in$R, f (x) > 0 positive (5.1) $\forall$x $\in$R, f (x) < 0 negative (5.2)

If (5.1) is formulated with the non-strict inequality f (x) $\geq$0, then we say that the function is non-negative. For the inequality f (x) $\leq$0, the term `non-positive' is uncommon, and one would normally use negative or zero. For example, the exponential function is positive, the absolute value function is non-negative, and the sine function is neither positive nor negative. Because inequalities are reversed under sign change, if a function f has any of the stated properties, then $-$f has the complementary property (e.g., if f is positive, then $-$f is negative).

© Springer-Verlag London 2014 F. Vivaldi, Mathematical Writing, Springer Undergraduate Mathematics Series, DOI 10.1007/978-1-4471-6527-9\_5

77

78 5 Describing Functions

Note that the terms `non-negative' and `not negative' have different meaning, the latter being the logical negation of negative. This distinction is very clear in the symbolic definitions.

$\forall$x $\in$R, f (x) $\geq$0 non-negative $\exists$x $\in$R, f (x) $\geq$0 not negative

(The second symbolic expression is obtained by negating (5.2) according to Theorem 4.3.) So a non-negative function is also not negative, but not vice-versa. Two similar expressions with different meanings can easily lead to confusion, and one must remain vigilant.

Next we consider how the action of a function affects the ordering of the real line; the order may be preserved, reversed, or a bit of both. There are two competing terminologies, labelled I and II in the table below. Each has advantages and disadvantages.

I II

$\forall$x, y $\in$R, x > y $\Rightarrow$f (x) > f (y) increasing strictly increasing $\forall$x, y $\in$R, x > y $\Rightarrow$f (x) $\geq$f (y) non-decreasing increasing $\forall$x, y $\in$R, x > y $\Rightarrow$f (x) < f (y) decreasing strictly decreasing $\forall$x, y $\in$R, x > y $\Rightarrow$f (x) $\leq$f (y) non-increasing decreasing

A function that is either increasing or decreasing (strictly or otherwise) is said to be monotonic.

Thus the arc-tangent is increasing in I and strictly increasing in II. In I no function can be both increasing and decreasing, and most functions are neither, for instance a constant or the cosine. The constants are the functions that are both non-decreasing and non-increasing.

The disadvantage of I is that, as we did above for `non-negative', we must differentiate between `non-increasing' and the logical negation of increasing ($\exists$x, y $\in$R, (x > y) $\land$( f (x) $\leq$f (y))).

The terminology II eliminates the annoying distinction between the prefixes `non-' and `not', but introduces a new problem. Now a constant is both increasing and decreasing, clashing with common usage. (After years without a pay rise, I wouldn't say: `my salary is increasing'.)

To be safe, use `strictly' to mean strict inequality in any case.

\section{5.2 Symmetries}

Real functions may have symmetries, expressing invariance with respect to changes of the argument. A function f is even or odd, respectively, if

$\forall$x $\in$R, f ($-$x) = f (x) or $\forall$x $\in$R, f ($-$x) = $-$f (x).

5.2 Symmetries 79

Fig. 5.1 An odd function

f(x)

x

So the cosine is even, the sine is odd, the exponential is neither even nor odd, and the zero function is both even and odd. The property of being even or odd has a geometrical meaning: graphs of even functions have a mirror symmetry with respect to the ordinate axis, while those of odd functions are symmetrical with respect to the origin (see Fig. 5.1). There is an easy way of constructing even/odd functions. For any real function g, the function x $\to$g(x) + g($-$x) is even, and so is the function g( f (x)) for any even function f . So the function x $\to$g(|x|) is even. To construct odd functions, we first define the sign function:

 

+1 if x > 0 0 if x = 0 $-$1 if x < 0

sign(x) =

x $\in$R. (5.3)



The sign function is odd: sign($-$x) = $-$sign(x). Then, for any real function g, the function x $\to$sign(x) g(|x|) is odd, as easily verified. This construct ensures that our function vanishes at the origin, which is a property of all odd functions.

Next we turn to translational symmetry, which is called periodicity. A function f is periodic with period T if

$\forall$x $\in$R, f (x + T ) = f (x) (5.4)

for some non-zero real number T (see Fig. 5.2).

For instance, the sine function is periodic with period T = 2$\pi$. If a function is periodic with period T , then it is also periodic with period 2T , 3T , etc. For this reason one normally requires the period T to be the smallest positive real number for which (5.4) is satisfied. To emphasise this point we use the terms minimal or fundamental period.

If we say that a function f is periodic---without reference to a specific period--- then the existence of the period must be required explicitly:

80 5 Describing Functions

Fig. 5.2 A periodic function

T f(x)

x

$\exists$T $\in$R+, $\forall$x $\in$R, f (x + T ) = f (x) (5.5)

where R+ is the set of positive real numbers---see (2.12). Note that the period T must be non-zero (lest this definition say nothing) and without loss of generality1

we shall require the period to be positive. [The case of negative period is dealt with the substitution x $\to$x $-$T in (5.4).]

The presence of symmetries reduces the amount of information needed to specify a function. If a function is even or odd, knowledge of the function for non-negative values of the argument suffices to characterise it completely. Likewise, if a periodic function is known over any interval of length equal to the period, then the function is specified completely.

One should be aware that symmetries are special; a function chosen `at random' will have no symmetry.

\section{5.3 Boundedness}

A set X $\subset$R is bounded if there is an interval containing it,2 namely if

$\exists$a, b $\in$R, $\forall$x $\in$X, a < x < b (5.6)

or

$\exists$c $\in$R, $\forall$x $\in$X, |x| < c. (5.7)

1 This expression indicates that an inessential restriction or simplification is being introduced. See Sect. 7.8 for another example. 2 The term `interval' is intended in the proper sense---infinite intervals are excluded, cf. (2.13).

5.3 Boundedness 81

The two definitions are equivalent. (Think about it.) The numbers a and b in (5.6) are an upper and a lower bound for X.

A real function f is bounded if its image f (R) is a bounded set. In symbols:

$\exists$c $\in$R, $\forall$x $\in$R, | f (x)| < c.

For example, the sine function is bounded and the exponential is not. The periodic function displayed in Fig. 5.2 is bounded.

A function f is bounded away from zero if its reciprocal is bounded. This means that for some positive constant c we have | f (x)| > c for all values of x. In symbols:

$\exists$c $\in$R+, $\forall$x $\in$R, | f (x)| > c.

The hyperbolic cosine is bounded away from zero (what could be a value of c in this case?) but the exponential function is not.

\section{5.4 Neighbourhoods}

A neighbourhood of a point x $\in$R is any open interval containing x. Although this definition makes no reference to the size of the interval, a neighbourhood of x contains all points sufficiently close to x. Thus the neighbourhood concept characterises `proximity' in a concise manner that does not require quantitative information. A skillful use of this term leads to terse and incisive statements. For instance, the sentence

The function f is bounded in a neighbourhood of x0.

means that there is an open interval containing x0 whose image under f is a bounded set. If we write this statement in symbols

$\exists$a, b $\in$R, $\exists$c $\in$R+, x0 $\in$(a, b) $\land$($\forall$x $\in$(a, b), | f (x)| < c)

we realise just how much information is packed into it. The following variant of the sentence above:

The function f is bounded in a sufficiently small neighbourhood of x0.

says exactly the same thing, but more eloquently; the reader is warned that the required neighbourhood may be very small.

A property of a function---boundedness in this case---which holds in a neighbourhood of some point of the domain of the function, but not necessarily in the whole domain, is said to be local. So a function may be locally increasing, locally injective, etc.

The function f has a maximum at x if the value of f at x is greater than the value at all other points, namely if

82 5 Describing Functions

$\forall$y $\in$R$\setminus$\{x\}, f (x) > f (y). (5.8)

The function f has a local maximum at x if (5.8) holds in some neighbourhood of x. The concept of minimum and local minimum are defined similarly. Thus the exponential has no maximum or minimum, the hyperbolic cosine has a minimum but no maximum, and the function x $\to$arctan(x)+sin(x) has infinitely many local maxima and minima, but no maximum or minimum.

By a neighbourhood of infinity we mean a ray \{x $\in$R : x > a\}, where a is a real number. A neighbourhood of $-$$\infty$is defined similarly, and both points at infinity $\pm$$\infty$are handled at once with the construct \{x $\in$R : |x| > a\}. The points $\pm$$\infty$are not numbers but they have neighbourhoods, which make them more tangible. So the sentence

The function f is constant in a neighbourhood of infinity.

means that f is constant for all sufficiently large values of the argument. It may be written symbolically as

$\exists$a $\in$R, $\forall$x $\in$R+, f (a + x) = f (a)

or as

$\exists$a $\in$R, $\forall$x > a, f (x) = f (a).

\section{5.4.1 Neighbourhoods and Sets}

Neighbourhoods are instrumental to the description of sets of numbers. This is an appealing part of the mathematical dictionary, due to the vivid mental pictures we associate with a geometric language. The case for expanding our dictionary is easily made:

X1 = \{1/n : n $\in$N\} X2 = Q $\cap$(0, 1).

How can we describe such sets?

Let X $\subset$R. A point x $\in$X is isolated if there is a neighbourhood of x that contains no other point of X. The set Z consists entirely of isolated points, and so does the set X1 above. By contrast, X2 has no isolated points.

A point x is an interior point of a set X if X contains a neighbourhood of x. A point x is a boundary point of a set X if every neighbourhood of x contains points of X as well as points of the complement of X. An isolated point is necessarily a boundary point, but not all boundary points are isolated. The boundary points of an interval are its end-points; all other points are interior points and there are no isolated points. Neither X1 nor X2 have interior points; the origin is a limit point of X1.

A set is closed if it contains all its boundary points, and is open if all its points are interior points. For intervals, these concepts agree with those given in Sect. 2.1.3.

5.4 Neighbourhoods 83

The sets X1 and X2 are neither open nor closed. The closure of a set X, denoted by

X, is the union of X and the boundary points of X. We see that X1 = X1 $\cup$\{0\} and X2 = [0, 1].

\section{5.5 Continuity}

Continuity is a fundamental concept in the theory of functions. Loosely speaking, a function is continuous if it has `no jumps', but this naive definition is only adequate for simple situations. We begin by considering continuity at a specific point of the domain of a function. An informal---yet accurate---characterisation of continuity is to say that a function is continuous at a point if its value there can be inferred unequivocally from the values at neighbouring points. Thus neighbourhoods come into play.

For example, the value of the sign function (5.3) at the origin cannot be inferred from the surrounding environment; even if we defined, say, sign(0) = 1, the ambiguity would remain. Things can get a lot worse. Consider the real function

1 if x $\in$Q

x $\to$

(5.9) 0 if x $\in$Q.

Given that in any neighbourhood of any real number there are both rational and irrational numbers, there is no way of inferring the value of this function at any point by considering how the function behaves in the surrounding region.

To write a symbolic definition of continuity we need to quantify neighbourhoods. For this purpose, let us denote by Nx the set of all neighbourhoods of a point x in R. (This is an abstract set, which may be represented as a set of end-points of real intervals---see Sect. 2.3.3 and Exercise 5.12.)

We say that f is continuous at the point a if for every neighbourhood J of f (a) we can find a neighbourhood of a whose image is contained in J. In symbols:

$\forall$J $\in$N f (a), $\exists$I $\in$Na, f (I) $\subset$J. (5.10)

This elaborate expression may be written in short-hand notation as

x$\to$a f (x) = f (a) lim

which says that f (a) is the limit of f (x) as x tends to a.

Let's go through Definition (5.10) in slow motion. The set Na consists of all open intervals containing the point a where continuity is being tested. The set N f (a) consists of all neighbourhoods of f (a), the value of the function at a. An arbitrary interval J containing f (a) is given, typically very small. We must now find an interval

I containing a whose image fits inside J. The choice of I will depend on J, and if

84 5 Describing Functions

we succeed in all cases, then the function is continuous at a. We don't need to know what the set f (I) looks like; it suffices to know that f (I) can be made small by making I small.

A function that is not continuous at a point of the domain is said to be discontinuous. A function is discontinuous everywhere if it has no points of continuity, like the function (5.9).

Continuity may be defined without reference to neighbourhoods, but in this case one must supply quantitative information. The full notation is considerably more involved than (5.10):

$\forall$$\varepsilon$ $\in$R+, $\exists$$\delta$ $\in$R+, $\forall$x $\in$R, |x $-$a| < $\delta$ $\Rightarrow$| f (x) $-$f (a)| < $\varepsilon$

and it remains so even if we strip it of all references to R:

$\forall$$\varepsilon$ > 0, $\exists$$\delta$ > 0, $\forall$x, |x $-$a| < $\delta$ $\Rightarrow$| f (x) $-$f (a)| < $\varepsilon$. (5.11)

A function is continuous if it is continuous at all points of the domain. An additional quantifier is required in Definition (5.10):

$\forall$a $\in$A, $\forall$J $\in$N f (a), $\exists$I $\in$Na, f (I) $\subset$J, (5.12)

where A is the domain of f .

A function f is differentiable at a if the limit

x$\to$a F(x) lim where F(x) = f (x) $-$f (a)

x = a x $-$a

exists. The function F is called the incremental ratio of f at a, which is not defined at x = a. However, if f is differentiable at a, then by letting F(a) def = limx$\to$a F(x), the function F becomes continuous at a. Let's sum this up in a sentence.

A function is differentiable at a point if its incremental ratio is continuous at that point; in this case the value of the incremental ratio is the derivative of the function.

A function is said to be differentiable if it is differentiable at all points of its domain. A function that is differentiable sufficiently often (all derivatives up to a sufficiently high order exist) is said to be smooth. The expression `sufficiently often' is deliberately vague; its precise meaning will depend on the context.

Many elementary real functions are continuous in their respective domains: the polynomials, the trigonometric functions, and the exponential and logarithmic functions, etc. These functions are also differentiable infinitely often.

A function f is singular at a point x if f is undefined at x, in which case x is a singularity (or a singular point) of the function. So the function x $\to$1/x is singular at the origin. A function that is not singular is said to be well-behaved or regular, the latter term used mostly for functions of a complex variable. We write:

5.5 Continuity 85

A rational function has finitely many singularities: the roots of the polynomial at the denominator. This function is regular everywhere else.

The tangent, secant, and co-secant have infinitely many singularities, evenly spaced along the real line.

In analysis the term `singular' is also used for well-defined functions that exhibit certain esoteric pathologies.

\section{5.6 Other Properties}

A function of the form f (x) = ax, where a is a real number, is said to be linear. The term linear originates from `line', and functions of the type f (x) = ax + b, are sometimes referred to as being `linear' because their graph is a line (the correct term is affine). This acceptation of the term linear is common for polynomial functions, which in the first instance are characterised by their degree. So we speak of a linear, quadratic, cubic, quartic function, etc.

Consider the absolute value function, defined as follows:

x if x $\geq$0

|x| =

(5.13) $-$x ifx < 0.

This function is not linear, but it is made of two linear pieces, glued together at the origin. Functions made of linear or affine pieces are said to be piecewise linear or piecewise affine [see Eq. (5.3) and Fig. 5.3]. A function is piecewise defined if its domain is partitioned into disjoint intervals or rays, with the function being specified independently over each interval. The properties of a piecewise-defined function may fail at the end-points of the intervals of definition. In this case terms such as piecewise increasing, piecewise continuous, piecewise differentiable, piecewise

Fig. 5.3 A continuous piecewise affine function

86 5 Describing Functions

Fig. 5.4 Left A step function. (The vertical segments are just a guide to the eye; they are not part of the graph of the function.) Right A periodic continuous function which is piecewise differentiable

smooth may be used---see Fig. 5.4. A piecewise constant function is called a step-function.

Example. The floor of a real number x, denoted by $\lfloor$x$\rfloor$, is the largest integer not exceeding x. Similarly, the ceiling $\lceil$x$\rceil$represents the smallest integer not smaller than x. Floor and ceiling are the prototype step functions. Closely connected to the floor is the fractional part of a real number x, denoted by \{x\}. (The notational clash with the set having x as its only element is one of the most spectacular in mathematics!) The fractional part is defined as

\{x\} := x $-$$\lfloor$x$\rfloor$

from which it follows that 0 $\leq$\{x\} < 1. This function is discontinuous and piecewise affine.

Example. Describe the following function: [ $\varepsilon$ ]

5.6 Other Properties 87

This is a smooth function, which is bounded and non-negative. It features an infinite sequence of evenly spaced local maxima, whose height decreases monotonically to zero for large arguments. There is one zero of the function between any two consecutive local maxima.

We rewrite it, borrowing some terminology from physics.

This function displays regular oscillations of constant period, with amplitude decreasing monotonically to zero.

The following functions have a behaviour that is qualitatively similar to that displayed above.

f (x) = sin2(x)e$-$x f (x) = 1 $-$cos(x)

. 1 + x2

\section{5.7 Describing Sequences}

A sequence may be thought of as a function defined over the natural numbers, or, more generally, over a subset of the integers (Sect. 3.1). Indeed in number theory the real (or complex) sequences are called arithmetical functions. Using this analogy, some terminology introduced for real functions translates literally to real sequences. So the terms

positive, negative, increasing, decreasing, monotonic, constant, periodic, bounded, bounded away from zero

have the same meaning for sequences as they have for functions. We write

The sequence of primes is positive, increasing, and unbounded.

Other terms require amendments or are simply not relevant to sequences. For instance, injectivity is not used---we just say that the terms of a sequence are distinct---while surjectivity is rarely significant. Invertibility is used in a different sense than for functions---see Sect. 9.4. The terms `even' and `odd' are still applicable to doubly-infinite sequences.

A sequence which settles down to a periodic pattern from a certain point on is said to be eventually periodic. More precisely, a sequence (x1, x2, . . .) is eventually periodic if for some k $\in$N the sub-sequence (xk, xk+1, . . .) is periodic. If the periodic pattern consists of a single term, then the sequence is said to be eventually constant.

The sequence of digits of a rational number is eventually periodic.

The idea of continuity does not apply to sequences because the concept of neighbourhood of an integer is pointless. Except at infinity, that is. The neighbourhoods of infinity are the infinite sets of the form \{k $\in$N : k > K\}, for some integer K, and we use the expression for all sufficiently large k, to mean for all k > K, for some K.

88 5 Describing Functions

We can now test continuity at infinity. Adapting Definition (5.10) to the present situation, we say that the sequence (ak) converges to the limit c if, given any neighbourhood J of c, all terms of the sequence eventually belong to J. In symbols:

$\forall$J $\in$Nc, $\exists$K $\in$N, $\forall$k > K, ak $\in$J.

This expression is written concisely as

k$\to$$\infty$ak = c. lim

By stipulating that a$\infty$= c, we see that a real sequence converges if it's continuous at infinity. In the special case c = $\infty$we say that the sequence diverges.

Exercise 5.1 Write each symbolic sentence in two ways:

(i) without any symbol, apart from f . (ii) with symbols only, using quantifiers. (You may assume that f : R $\to$R.)

1. f (0) $\in$Q 2. f (R) = R

3. \# f (R) = 1 4. f (Z) = \{0\}

5. 0 $\in$f (Z) 6. f $-$1(\{0\}) = Z

7. f (R) $\subset$Q 8. f (R) $\supset$Z

9. f (Z) = f (N) 10. f (Q) $\cap$Q = $\emptyset$

11. f $-$1(Q) = $\emptyset$ 12. \# f $-$1(Z) < $\infty$.

Exercise 5.2 For each symbolic sentence of the previous exercise, give examples of functions for which the sentence is true and functions for which the sentence is false.

Exercise 5.3 Write each symbolic sentence without symbols, apart from f .

1. $\forall$x $\in$Z, f (2x) = 0

2. $\forall$x $\in$R, f ( f (x)) = x

3. $\forall$x $\in$R+, f ($-$x) = 0

4. $\forall$x $\in$[0, 1], f (x) = 0

5. $\forall$x $\in$Q$\setminus$\{0\}, f (x) = 0

6. $\forall$x $\in$N, f (x) $\in$Q

7. $\forall$x $\in$Z, $\exists$y $\in$N, f (x + y) = 0

8. $\forall$x $\in$R, $\exists$y $\in$R, | f (y)| < | f (x)|/2

9. $\forall$x $\in$R, f (x) = 4 f (x/2).

5.7 Describing Sequences 89

Exercise 5.4 Turn words into symbols.

1. The function f has no zero.

2. The function f is not always positive.

3. The function f does not vanish except, possibly, at the origin.

4. The function f vanishes identically outside the open unit interval.

5. The zeros of f include all multiples of $\pi$.

6. The function f may vanish only at some rational.

7. There are zeros of f arbitrarily close to the origin.

8. The zeros of f are isolated.

9. There is a zero of f between any two consecutive integers.

10. The function f has no zero for all sufficiently large arguments.

11. Every integer is a local maximum of f .

12. The function f is locally invertible near the origin.

Exercise 5.5 Consider the following implications, where f is a real function.

1. If f is decreasing, then $-$f is increasing.

2. If f is decreasing, then | f | is increasing.

3. If | f | is increasing, then f is monotonic.

4. If f is even, then | f | is also even.

5. If f is unbounded, then f is surjective.

6. If f is continuous, then | f | is continuous.

7. If f is differentiable, then | f | is differentiable.

8. If f is periodic, then f 2 is periodic.

Of each implication:

(a) state the converse, and decide whether it's true or false; (b) state the contrapositive, and decide whether it's true or false;

(c) state the negation.

Exercise 5.6 As in Exercise 4.3.

1. Every differentiable real function is continuous.

2. A bounded real function cannot be injective.

3. The sum of two odd functions is an odd function.

4. No polynomial function is bounded.

5. The composition of two periodic functions is periodic.

6. Any function with an increasing derivative is increasing.

Exercise 5.7 Consider the Definition (5.5) of a periodic function. What happens if we place the quantifiers in reverse order?

$\forall$x $\in$R, $\exists$T $\in$R$*$, f (x + T ) = f (x) (5.14)

$\forall$x $\in$R, $\exists$T $\in$R+, f (x + T ) = f (x). (5.15)

90 5 Describing Functions

1. Find a function that satisfies (5.14) but not (5.5).

2. Find a function that satisfies (5.15) but not (5.5).

3. Characterize the set of functions that satisfy (5.14). Do the same for (5.15).

Exercise 5.8 Describe the behaviour of the following functions. [ $\varepsilon$, 30]

1) 2)

3) 4)

y

5) 6)

x x

5.7 Describing Sequences 91

8)

7)

x

9) 10)

Exercise 5.9 By experimenting on a computer, if necessary, define real functions whose behaviour is qualitatively similar to those of the previous problem.

Exercise 5.10 Give an example of a real function with the stated properties.

1. Increasing, with decreasing derivative.

2. Unbounded, with unbounded derivative.

3. Unbounded, with bounded derivative.

4. Discontinuous, with continuous absolute value.

5. With infinitely many zeros, but not periodic.

6. Unbounded and vanishing infinitely often.

7. The image is a half-open interval.

8. The image is the set of natural numbers.

Exercise 5.11 Express each statement with symbols.

1. The sequences (ak) and (bk) are distinct.

2. The sequence (ak) is eventually constant.

3. The sequence (ak) is not periodic.

4. The sequence (ak) is eventually periodic.

5. The sequence (ak) has infinitely many negative terms.

6. Eventually, all terms of the sequence (ak) become negative.

92 5 Describing Functions

7. The terms of the sequence (ak) get arbitrarily close to zero.

8. The sequence (ak) does not converge.

9. Each term of the sequence (ak) appears at least twice.

10. Each term of the sequence (ak) appears infinitely often.

11. Each natural number appears infinitely often in the sequence (ak).

12. The sequence (ak) has a bounded sub-sequence.

13. The sequence (ak) has a convergent sub-sequence.

Exercise 5.12 Find a representation for the abstract set Nx of all neighbourhoods of a point x $\in$R (see Sect. 2.3.3).

\end{document}